```latex
\documentclass[12pt,a4paper]{article}

\usepackage[utf8]{inputenc}
\usepackage[T1]{fontenc}
\usepackage{geometry}
\usepackage{enumitem}
\usepackage{longtable}
\usepackage{array}
\usepackage{booktabs}
\usepackage{hyperref}
\usepackage{xcolor}
\usepackage{listings}
\usepackage{tikz}
\usepackage{float}

\geometry{margin=1in}

\title{
    \textbf{ATM Management and Banking Transaction Simulation System}\\
    \large MCA Final Year Project Specification
}

\author{MCA Final Project}
\date{\today}

\begin{document}

\maketitle

\tableofcontents
\newpage


% ============================================================
\section{Project Overview}
% ============================================================

The project is a terminal-based ATM Management and Banking Transaction
Simulation System developed entirely in the C programming language.

The system will simulate the major operations performed by a real ATM,
including card authentication, PIN verification, balance enquiry, cash
withdrawal, cash deposit, fund transfer, mini statement, PIN change,
transaction management, ATM cash management, and administrative operations.

The application will use a MySQL relational database for persistent data
storage.

The project must be developed as a modular C application and must not use
C++, C\#, Java, Python, PHP, JavaScript, or a web frontend.


% ============================================================
\section{Technology Stack}
% ============================================================

\begin{longtable}{p{0.3\textwidth} p{0.6\textwidth}}
\toprule
\textbf{Component} & \textbf{Technology} \\
\midrule
Programming Language & C \\
User Interface & Terminal / Command Line Interface (CLI) \\
Database & MySQL \\
Database Connectivity & MySQL C API / MySQL Connector for C \\
Operating System & Windows / Linux compatible where practical \\
Compiler & GCC or compatible C compiler \\
Version Control & Git (optional but recommended) \\
\bottomrule
\end{longtable}


% ============================================================
\section{Project Objectives}
% ============================================================

The primary objectives are:

\begin{enumerate}
    \item Develop a realistic ATM simulation using C.
    \item Demonstrate modular programming in C.
    \item Demonstrate relational database concepts using MySQL.
    \item Implement secure card and PIN authentication.
    \item Implement banking transaction operations.
    \item Maintain persistent customer and transaction data.
    \item Demonstrate database transactions and rollback.
    \item Implement ATM cash inventory management.
    \item Implement administrative management functions.
    \item Provide proper input validation and error handling.
    \item Generate transaction records and mini statements.
    \item Produce an academically presentable MCA final-year project.
\end{enumerate}


% ============================================================
\section{Project Scope}
% ============================================================

The system will contain two primary roles:

\begin{enumerate}
    \item Customer
    \item Administrator
\end{enumerate}

\subsection{Customer Scope}

Customers will be able to:

\begin{itemize}
    \item Insert/use a registered card.
    \item Authenticate using a PIN.
    \item View account balance.
    \item Withdraw cash.
    \item Deposit cash.
    \item Transfer funds.
    \item View mini statement.
    \item Change PIN.
    \item Logout safely.
\end{itemize}

\subsection{Administrator Scope}

Administrators will be able to:

\begin{itemize}
    \item Login to the administration panel.
    \item View customers.
    \item View accounts.
    \item Search accounts.
    \item Create customer/account records.
    \item Block or unblock cards.
    \item View transactions.
    \item View ATM cash inventory.
    \item Refill ATM cash.
    \item View transaction statistics.
\end{itemize}


% ============================================================
\section{High-Level System Architecture}
% ============================================================

The application will follow a modular layered architecture.

\begin{verbatim}
                    ATM SIMULATION SYSTEM
                            |
                    C Terminal Application
                            |
          +-----------------+------------------+
          |                 |                  |
     Authentication      Customer           Admin
          Module         Operations         Operations
          |                 |                  |
          +-----------------+------------------+
                            |
                     Business Logic
                            |
                     Database Layer
                            |
                         MySQL
\end{verbatim}

The system should separate:

\begin{enumerate}
    \item User Interface logic
    \item Authentication logic
    \item Business logic
    \item Database operations
    \item Utility functions
\end{enumerate}

Database queries should not be scattered throughout every source file.


% ============================================================
\section{Recommended Project Structure}
% ============================================================

The project should use the following structure:

\begin{verbatim}
ATM-Simulation/
|
+-- src/
|   +-- main.c
|   +-- auth.c
|   +-- customer.c
|   +-- account.c
|   +-- card.c
|   +-- transaction.c
|   +-- withdrawal.c
|   +-- deposit.c
|   +-- transfer.c
|   +-- atm.c
|   +-- admin.c
|   +-- database.c
|   +-- security.c
|   +-- ui.c
|   +-- receipt.c
|   +-- validation.c
|   +-- utils.c
|
+-- include/
|   +-- auth.h
|   +-- customer.h
|   +-- account.h
|   +-- card.h
|   +-- transaction.h
|   +-- withdrawal.h
|   +-- deposit.h
|   +-- transfer.h
|   +-- atm.h
|   +-- admin.h
|   +-- database.h
|   +-- security.h
|   +-- ui.h
|   +-- receipt.h
|   +-- validation.h
|   +-- utils.h
|
+-- database/
|   +-- schema.sql
|   +-- seed.sql
|
+-- receipts/
|
+-- tests/
|
+-- docs/
|
+-- Makefile
|
+-- README.md
\end{verbatim}


% ============================================================
\section{Database Design}
% ============================================================

The system will use MySQL as the relational database management system.

The database name should be:

\begin{verbatim}
atm_simulation
\end{verbatim}

The database should contain the following major tables:

\begin{enumerate}
    \item customers
    \item accounts
    \item cards
    \item beneficiaries
    \item transactions
    \item atm
    \item atm\_cash
    \item admins
\end{enumerate}


% ============================================================
\subsection{Customers Table}
% ============================================================

\begin{verbatim}
customers
---------
customer_id       PRIMARY KEY
full_name
mobile
email
address
created_at
status
\end{verbatim}

Possible status values:

\begin{verbatim}
ACTIVE
INACTIVE
\end{verbatim}


% ============================================================
\subsection{Accounts Table}
% ============================================================

\begin{verbatim}
accounts
--------
account_id        PRIMARY KEY
customer_id       FOREIGN KEY
account_number    UNIQUE
account_type
balance
daily_limit
account_status
created_at
\end{verbatim}

Possible account types:

\begin{verbatim}
SAVINGS
CURRENT
\end{verbatim}

Possible account statuses:

\begin{verbatim}
ACTIVE
BLOCKED
CLOSED
\end{verbatim}


% ============================================================
\subsection{Cards Table}
% ============================================================

\begin{verbatim}
cards
-----
card_id
account_id
card_number       UNIQUE
pin_hash
failed_attempts
card_status
issue_date
expiry_date
\end{verbatim}

Possible card statuses:

\begin{verbatim}
ACTIVE
BLOCKED
EXPIRED
CANCELLED
\end{verbatim}

The PIN must not be stored as plain text.


% ============================================================
\subsection{Beneficiaries Table}
% ============================================================

\begin{verbatim}
beneficiaries
-------------
beneficiary_id
account_id
beneficiary_account_number
beneficiary_name
nickname
status
created_at
\end{verbatim}


% ============================================================
\subsection{Transactions Table}
% ============================================================

\begin{verbatim}
transactions
------------
transaction_id
transaction_reference
account_id
transaction_type
amount
balance_before
balance_after
description
status
transaction_date
\end{verbatim}

Transaction types may include:

\begin{verbatim}
WITHDRAWAL
DEPOSIT
TRANSFER
BALANCE_ENQUIRY
PIN_CHANGE
\end{verbatim}

Transaction status:

\begin{verbatim}
SUCCESS
FAILED
CANCELLED
\end{verbatim}


% ============================================================
\subsection{ATM Table}
% ============================================================

\begin{verbatim}
atm
---
atm_id
atm_code
location
status
created_at
\end{verbatim}


% ============================================================
\subsection{ATM Cash Table}
% ============================================================

The ATM should maintain physical cash inventory by denomination.

\begin{verbatim}
atm_cash
--------
cash_id
atm_id
denomination
quantity
updated_at
\end{verbatim}

Example:

\begin{verbatim}
Denomination     Quantity

Rs. 500             120
Rs. 200              80
Rs. 100             100
Rs. 50               40
\end{verbatim}


% ============================================================
\subsection{Admins Table}
% ============================================================

\begin{verbatim}
admins
------
admin_id
username
password_hash
status
created_at
\end{verbatim}


% ============================================================
\section{Database Relationships}
% ============================================================

The primary relationships are:

\begin{verbatim}
CUSTOMER
   |
   | 1:N
   v
ACCOUNT
   |
   | 1:N
   +----------> CARD
   |
   +----------> TRANSACTION
   |
   +----------> BENEFICIARY

ATM
 |
 +----------> ATM_CASH

ADMIN
 |
 +----------> ADMIN OPERATIONS
\end{verbatim}


% ============================================================
\section{Customer Authentication Flow}
% ============================================================

The customer authentication process must be:

\begin{verbatim}
START
  |
  v
INSERT CARD
  |
  v
ENTER CARD NUMBER
  |
  v
CARD EXISTS?
  |
  +---- NO ----> INVALID CARD
  |
 YES
  |
  v
CHECK CARD STATUS
  |
  +---- BLOCKED/EXPIRED/CANCELLED
  |              |
  |              v
  |         DENY ACCESS
  |
 ACTIVE
  |
  v
ENTER PIN
  |
  v
VERIFY PIN HASH
  |
  +---- WRONG ----> INCREMENT FAILED ATTEMPTS
  |                      |
  |                      +---- 3 attempts ----> BLOCK CARD
  |
 CORRECT
  |
  v
RESET FAILED ATTEMPTS
  |
  v
CREATE SESSION
  |
  v
CUSTOMER MAIN MENU
\end{verbatim}


% ============================================================
\section{PIN Security Rules}
% ============================================================

The following rules are mandatory:

\begin{enumerate}
    \item Never store a PIN in plain text.
    \item Store a secure representation of the PIN.
    \item Do not display the entered PIN on screen.
    \item Maximum 3 failed PIN attempts.
    \item Block the card after 3 consecutive failed attempts.
    \item Reset failed attempts after successful authentication.
    \item PIN change requires verification of the current PIN.
\end{enumerate}


% ============================================================
\section{Customer Main Menu}
% ============================================================

After successful authentication:

\begin{verbatim}
========================================
             ATM MAIN MENU
========================================

1. Balance Enquiry
2. Cash Withdrawal
3. Cash Deposit
4. Fund Transfer
5. Mini Statement
6. Change PIN
7. Logout

Enter Choice:
\end{verbatim}


% ============================================================
\section{Balance Enquiry}
% ============================================================

The system should:

\begin{enumerate}
    \item Identify the authenticated account.
    \item Retrieve current balance.
    \item Display balance.
    \item Optionally create a balance enquiry transaction record.
\end{enumerate}


% ============================================================
\section{Cash Withdrawal}
% ============================================================

Withdrawal must validate:

\begin{enumerate}
    \item Account status.
    \item Valid withdrawal amount.
    \item Minimum withdrawal amount.
    \item Maximum withdrawal amount.
    \item Daily withdrawal limit.
    \item Sufficient account balance.
    \item Sufficient ATM cash.
    \item Available denominations.
\end{enumerate}

The withdrawal operation should use a database transaction.

\begin{verbatim}
BEGIN TRANSACTION

1. Validate balance
2. Validate ATM cash
3. Calculate denomination distribution
4. Update account balance
5. Update ATM cash inventory
6. Insert transaction record

COMMIT
\end{verbatim}

If any step fails:

\begin{verbatim}
ROLLBACK
\end{verbatim}


% ============================================================
\section{ATM Denomination Logic}
% ============================================================

The ATM should not simply reduce the total cash amount.

It should calculate actual available notes.

Example:

For a withdrawal of Rs. 2700:

\begin{verbatim}
Rs. 500 x 5 = Rs. 2500
Rs. 100 x 2 = Rs.  200
-----------------------
Total       = Rs. 2700
\end{verbatim}

The algorithm must ensure that the required combination is actually
available in the ATM.

If the exact amount cannot be formed from available denominations,
the withdrawal should be rejected rather than incorrectly modifying
the account balance.


% ============================================================
\section{Cash Deposit}
% ============================================================

Deposit flow:

\begin{verbatim}
ENTER AMOUNT
     |
VALIDATE AMOUNT
     |
BEGIN TRANSACTION
     |
UPDATE ACCOUNT BALANCE
     |
UPDATE ATM CASH
     |
INSERT TRANSACTION
     |
COMMIT
\end{verbatim}

The system must reject:

\begin{verbatim}
0
negative values
non-numeric values
unreasonably large values
\end{verbatim}


% ============================================================
\section{Fund Transfer}
% ============================================================

The customer should be able to transfer money to another valid account.

Flow:

\begin{verbatim}
ENTER BENEFICIARY / ACCOUNT NUMBER
              |
        VALIDATE ACCOUNT
              |
        ENTER TRANSFER AMOUNT
              |
       CHECK AVAILABLE BALANCE
              |
       BEGIN TRANSACTION
              |
       DEBIT SENDER ACCOUNT
              |
       CREDIT RECEIVER ACCOUNT
              |
       INSERT TRANSACTION RECORDS
              |
            COMMIT
\end{verbatim}

If any step fails:

\begin{verbatim}
ROLLBACK
\end{verbatim}

The sender and receiver balances must never become inconsistent.


% ============================================================
\section{Mini Statement}
% ============================================================

The system should display the latest 5 or 10 transactions.

Example:

\begin{verbatim}
---------------------------------------------------------
DATE          TYPE             AMOUNT        BALANCE
---------------------------------------------------------
07-10-2026    WITHDRAWAL       Rs. 5000      Rs. 20000
07-10-2026    DEPOSIT          Rs. 3000      Rs. 25000
06-10-2026    TRANSFER         Rs. 1500      Rs. 22000
05-10-2026    WITHDRAWAL       Rs. 2000      Rs. 23500
04-10-2026    DEPOSIT          Rs. 5000      Rs. 25500
---------------------------------------------------------
\end{verbatim}


% ============================================================
\section{PIN Change}
% ============================================================

PIN change process:

\begin{verbatim}
ENTER CURRENT PIN
       |
VERIFY CURRENT PIN
       |
ENTER NEW PIN
       |
CONFIRM NEW PIN
       |
CHECK NEW PIN VALIDITY
       |
STORE NEW PIN HASH
       |
SUCCESS
\end{verbatim}

The system must reject:

\begin{itemize}
    \item Incorrect current PIN.
    \item New PIN and confirmation mismatch.
    \item Invalid PIN length.
    \item Weak PIN patterns where appropriate.
\end{itemize}


% ============================================================
\section{Session Management}
% ============================================================

Once the user successfully authenticates, the application should maintain
a session containing the authenticated account/card information.

The user must not be asked for the card number again for every transaction.

The session ends when:

\begin{enumerate}
    \item The user selects Logout.
    \item The session expires due to inactivity.
    \item A critical authentication/security event occurs.
\end{enumerate}


% ============================================================
\section{Card States}
% ============================================================

The system should support:

\begin{verbatim}
ACTIVE
BLOCKED
EXPIRED
CANCELLED
\end{verbatim}

Blocked, expired, or cancelled cards must not be allowed to perform
transactions.


% ============================================================
\section{Transaction IDs}
% ============================================================

Every financial transaction should receive a unique transaction
reference.

Example:

\begin{verbatim}
TXN20261007142301
\end{verbatim}

Transaction references must be unique.


% ============================================================
\section{Receipt Generation}
% ============================================================

After successful financial transactions, the system should optionally
generate a text receipt.

Example:

\begin{verbatim}
========================================
              XYZ BANK ATM
========================================

Transaction ID : TXN20261007142301
Card Number    : XXXX XXXX XXXX 1234
Transaction    : CASH WITHDRAWAL
Amount         : Rs. 5,000
Balance        : Rs. 20,000

Date           : 07-10-2026
Time           : 14:23:01
Status         : SUCCESS

        THANK YOU FOR BANKING
========================================
\end{verbatim}

Receipts may be stored inside the \texttt{receipts/} directory.


% ============================================================
\section{Admin Panel}
% ============================================================

The administrator must authenticate before accessing administrative
operations.

Admin menu:

\begin{verbatim}
========================================
             ADMIN PANEL
========================================

1. View Customers
2. Search Customer
3. View Accounts
4. Create Customer
5. Create Account
6. Block Card
7. Unblock Card
8. View Transactions
9. ATM Cash Status
10. Refill ATM
11. Transaction Statistics
12. Logout

Enter Choice:
\end{verbatim}


% ============================================================
\section{ATM Cash Management}
% ============================================================

The administrator should be able to view:

\begin{verbatim}
ATM CASH INVENTORY
-----------------------------
Rs. 500 : 120 notes
Rs. 200 :  80 notes
Rs. 100 : 100 notes
Rs.  50 :  40 notes
-----------------------------
TOTAL   : Rs. 88,000
\end{verbatim}

The administrator should be able to refill the ATM.

The system must never allow the ATM cash balance to become negative.


% ============================================================
\section{Transaction Statistics}
% ============================================================

The admin panel may display:

\begin{verbatim}
Today's Statistics
------------------------------
Total Transactions : 127
Withdrawals        : 62
Deposits           : 31
Transfers          : 28
Failed             : 6
------------------------------
ATM Cash           : Rs. 124500
\end{verbatim}


% ============================================================
\section{Input Validation}
% ============================================================

Every user input must be validated.

The program must safely handle:

\begin{itemize}
    \item Empty input.
    \item Non-numeric input.
    \item Negative numbers.
    \item Zero values.
    \item Very large values.
    \item Invalid account numbers.
    \item Invalid card numbers.
    \item Invalid menu selections.
\end{itemize}

The program must not crash because of normal invalid user input.


% ============================================================
\section{Error Handling}
% ============================================================

The application must display clear error messages.

Examples:

\begin{verbatim}
ERROR: Invalid card number.

ERROR: Incorrect PIN.
Attempts remaining: 2

ERROR: Card is blocked.

ERROR: Insufficient account balance.

ERROR: ATM does not have sufficient cash.

ERROR: Requested denomination combination is unavailable.

ERROR: Beneficiary account does not exist.

ERROR: Database operation failed.
\end{verbatim}


% ============================================================
\section{Database Transaction and ACID Requirements}
% ============================================================

Financial operations must use database transactions where multiple
database modifications must remain consistent.

Examples include:

\begin{itemize}
    \item Cash withdrawal.
    \item Cash deposit.
    \item Fund transfer.
    \item ATM cash update.
\end{itemize}

The implementation should demonstrate:

\begin{verbatim}
BEGIN
COMMIT
ROLLBACK
\end{verbatim}

The system must avoid situations where:

\begin{verbatim}
Account balance updated
        +
ATM cash not updated
```

or:

\begin{verbatim}
Sender debited
        +
Receiver not credited
```

These operations must be atomic.


% ============================================================
\section{Security Requirements}
% ============================================================

The following security practices should be implemented:

\begin{enumerate}
    \item PIN should not be stored in plain text.
    \item Password/PIN input should be masked where practical.
    \item SQL queries should use parameterized/prepared statements where
          supported by the C MySQL API.
    \item User input must not be directly concatenated into SQL queries.
    \item Card status must be checked before transactions.
    \item Account status must be checked before transactions.
    \item Failed PIN attempts must be tracked.
    \item Sensitive information should not be printed unnecessarily.
\end{enumerate}


% ============================================================
\section{Terminal User Interface}
% ============================================================

The terminal interface should look professional despite being CLI based.

Use:

\begin{itemize}
    \item Clear headings.
    \item Consistent separators.
    \item Numbered menus.
    \item Success messages.
    \item Error messages.
    \item Confirmation messages.
    \item Screen clearing where appropriate.
    \item Safe cross-platform terminal behavior where practical.
\end{itemize}

Avoid depending heavily on emojis because terminal encoding can differ
between systems.

Example:

\begin{verbatim}
+--------------------------------------+
|          ATM SIMULATION              |
+--------------------------------------+
|                                      |
|        Welcome to XYZ Bank           |
|                                      |
+--------------------------------------+
\end{verbatim}


% ============================================================
\section{Application Flow}
% ============================================================

Main application:

\begin{verbatim}
START
 |
 v
DATABASE CONNECTION
 |
 +---- FAILURE ----> DISPLAY ERROR AND EXIT
 |
 SUCCESS
 |
 v
WELCOME SCREEN
 |
 +---------------------------+
 |                           |
Customer                   Admin
 |                           |
Card Input                Admin Login
 |                           |
PIN Verification          Admin Dashboard
 |                           |
Customer Menu             Admin Operations
 |
Operations
 |
Logout
 |
END
\end{verbatim}


% ============================================================
\section{Important Business Rules}
% ============================================================

The following rules must be enforced.

\begin{enumerate}
    \item Only registered cards can access customer operations.
    \item Blocked, expired, and cancelled cards cannot authenticate.
    \item Maximum three consecutive incorrect PIN attempts.
    \item Successful authentication resets failed attempt count.
    \item Withdrawal cannot exceed account balance.
    \item Withdrawal cannot exceed the daily withdrawal limit.
    \item Withdrawal cannot exceed available ATM cash.
    \item ATM must have the required denominations.
    \item Account balance cannot become negative.
    \item Transfer requires a valid receiver account.
    \item Sender must have sufficient balance.
    \item Financial transactions must be atomic.
    \item Every successful financial transaction must be recorded.
    \item Every transaction must have a unique transaction reference.
    \item Logout must terminate the authenticated session.
    \item Invalid input must not crash the application.
\end{enumerate}


% ============================================================
\section{Testing Requirements}
% ============================================================

The project must include test cases for at least:

\begin{longtable}{p{0.08\textwidth} p{0.42\textwidth} p{0.35\textwidth}}
\toprule
\textbf{ID} & \textbf{Test Case} & \textbf{Expected Result} \\
\midrule
T01 & Valid card + valid PIN & Login successful \\
T02 & Invalid card & Access denied \\
T03 & Wrong PIN once & Error + remaining attempts \\
T04 & Wrong PIN three times & Card blocked \\
T05 & Blocked card login & Access denied \\
T06 & Balance enquiry & Correct balance displayed \\
T07 & Valid withdrawal & Balance and ATM cash updated \\
T08 & Insufficient balance & Withdrawal rejected \\
T09 & Insufficient ATM cash & Withdrawal rejected \\
T10 & Invalid withdrawal amount & Validation error \\
T11 & Valid deposit & Balance and ATM cash updated \\
T12 & Invalid transfer account & Transfer rejected \\
T13 & Valid transfer & Sender and receiver updated \\
T14 & Mini statement & Recent transactions displayed \\
T15 & Correct PIN change & PIN updated \\
T16 & Wrong current PIN & PIN change rejected \\
T17 & Database failure & Safe error handling \\
T18 & Transaction failure & Database rollback \\
T19 & Admin login & Admin dashboard displayed \\
T20 & ATM refill & Cash inventory updated \\
\bottomrule
\end{longtable}


% ============================================================
\section{Non-Functional Requirements}
% ============================================================

\begin{enumerate}
    \item The application should be modular.
    \item The application should be readable and maintainable.
    \item Database operations should be centralized.
    \item Functions should have clear responsibilities.
    \item The application should not crash on normal invalid input.
    \item Error messages should be understandable.
    \item SQL queries should be safe.
    \item Financial operations should maintain data consistency.
\end{enumerate}


% ============================================================
\section{Development Phases}
% ============================================================

Development should be performed in the following order.

\subsection{Phase 1 -- Environment Setup}

\begin{itemize}
    \item Install/configure C compiler.
    \item Install MySQL.
    \item Install/configure MySQL C Connector.
    \item Verify C-to-MySQL connection.
\end{itemize}

\subsection{Phase 2 -- Database}

\begin{itemize}
    \item Create database.
    \item Create tables.
    \item Add primary keys.
    \item Add foreign keys.
    \item Add constraints.
    \item Add indexes.
    \item Insert demo data.
\end{itemize}

\subsection{Phase 3 -- C Foundation}

\begin{itemize}
    \item Create project structure.
    \item Implement database connection module.
    \item Implement UI utilities.
    \item Implement input validation.
    \item Implement common utilities.
\end{itemize}

\subsection{Phase 4 -- Authentication}

\begin{itemize}
    \item Card lookup.
    \item Card status validation.
    \item PIN verification.
    \item Failed attempt tracking.
    \item Card blocking.
    \item Session creation.
\end{itemize}

\subsection{Phase 5 -- Customer Operations}

\begin{itemize}
    \item Balance enquiry.
    \item Withdrawal.
    \item Deposit.
    \item Transfer.
    \item Mini statement.
    \item PIN change.
    \item Logout.
\end{itemize}

\subsection{Phase 6 -- ATM Management}

\begin{itemize}
    \item Denomination management.
    \item ATM cash calculation.
    \item Withdrawal denomination algorithm.
    \item ATM refill.
\end{itemize}

\subsection{Phase 7 -- Admin Panel}

\begin{itemize}
    \item Admin authentication.
    \item Customer management.
    \item Account management.
    \item Card management.
    \item Transaction reports.
    \item ATM cash management.
\end{itemize}

\subsection{Phase 8 -- Testing}

Perform unit testing, integration testing, database testing and
transaction rollback testing.

\subsection{Phase 9 -- Documentation}

Prepare complete MCA project documentation.


% ============================================================
\section{Expected Final Deliverable}
% ============================================================

The final system should provide:

\begin{itemize}
    \item A working C terminal application.
    \item MySQL database.
    \item Database schema SQL file.
    \item Demo/seed SQL data.
    \item Modular C source code.
    \item Header files.
    \item Build instructions.
    \item Test cases.
    \item Sample transaction receipts.
    \item Project documentation.
    \item ER diagram.
    \item DFD.
    \item Flowcharts.
    \item Database relational schema.
\end{itemize}


% ============================================================
\section{Important Development Rule for AI-Assisted Development}
% ============================================================

If an AI coding assistant such as Antigravity is used, it must not
generate the entire project blindly in a single step.

The project must be developed incrementally.

The AI must:

\begin{enumerate}
    \item Understand the architecture before writing code.
    \item Confirm dependencies before using them.
    \item Keep modules separated.
    \item Avoid duplicate functions.
    \item Avoid global mutable state unless necessary.
    \item Use the database schema as the source of truth.
    \item Compile the project after meaningful changes.
    \item Fix compiler warnings.
    \item Test database operations.
    \item Test financial transaction rollback.
    \item Never silently change business rules.
\end{enumerate}


% ============================================================
\section{Future Scope}
% ============================================================

Possible future enhancements include:

\begin{itemize}
    \item Multiple ATM locations.
    \item Bank-to-bank transfers.
    \item OTP-based authentication.
    \item Account statement export.
    \item PDF receipt generation.
    \item GUI frontend.
    \item Role-based administrative permissions.
    \item Audit logs.
    \item Advanced fraud detection.
    \item Network-based ATM server architecture.
\end{itemize}


% ============================================================
\section{Project Philosophy}
% ============================================================

The objective is not to create a large amount of code.

The objective is to create a logically correct, modular, database-driven
ATM simulation that demonstrates:

\begin{verbatim}
C Programming
      +
Data Structures
      +
Functions and Modules
      +
SQL
      +
Relational Database
      +
Authentication
      +
Transaction Management
      +
Error Handling
      +
Security
      +
Software Engineering
\end{verbatim}

The final application should therefore prioritize correctness,
maintainability, database consistency and realistic ATM behavior over
unnecessary features.

\end{document}
```